\section{Introducción}

El presente trabajo desarrolla la medición de calidad en uso para el caso MediSalud HIS mediante ISO/IEC 25022. El sistema hospitalario comprende procesos de agenda, admisión, historia clínica electrónica, enfermería, farmacia, facturación, telemedicina y atención al paciente en cinco sedes. La evaluación utiliza 3.000 incidentes registrados durante 2025 y evidencia simulada para observar la efectividad, eficiencia, satisfacción, libertad de riesgo y cobertura del contexto.

La aplicación local recrea tareas clínicas, administrativas y del paciente con fallos controlados. Su propósito es producir registros reproducibles para el taller; no procesa información clínica real ni representa un despliegue productivo. Los datos originales se conservan sin modificaciones y se distinguen de los eventos, encuestas y resultados derivados.

\subsection{Objetivo general}

Evaluar la calidad en uso de MediSalud HIS mediante ISO/IEC 25022, relacionando usuarios, tareas, contextos, incidentes y métricas obtenidas a partir del caso académico y de una simulación local reproducible.

\subsection{Objetivos específicos}

\begin{enumerate}
  \item Analizar los procesos, perfiles y problemas descritos en el caso MediSalud.
  \item Clasificar los 3.000 incidentes según las características de calidad en uso.
  \item Diferenciar calidad interna, externa y en uso dentro del modelo SQuaRE.
  \item Definir atributos y métricas con fórmula, unidad, fuente, meta y semáforo.
  \item Generar eventos y encuestas simuladas con semilla y procedencia identificables.
  \item Automatizar el cálculo de KPI e interpretar los resultados obtenidos.
  \item Formular un plan de mejora posterior sin alterar la línea base evaluada.
\end{enumerate}
